% ============================================================================
% Beyond the static gradient: electrons on helium as a time-dependent-field
% sensor for particle physics
% Target: PRX Quantum / Phys. Rev. A (adjust class options below once decided)
% ============================================================================
\documentclass[aps,prxquantum,reprint,superscriptaddress,longbibliography]{revtex4-2}
% \documentclass[aps,pra,reprint,superscriptaddress,longbibliography]{revtex4-2}          % Phys. Rev. A

\usepackage{graphicx}
\usepackage{amsmath,amssymb}
\usepackage{bm}
\usepackage{hyperref}

% Shared project conventions (keep consistent with
% skills/literature-and-notation/references/notation.md):
%   natural units hbar^2/2m_e = 1, hbar = 1
\newcommand{\ehe}{electrons on helium}

\begin{document}

\title{Beyond the Static Gradient: Electrons on Helium as a Time-Dependent-Field Sensor for
Particle Physics}
% Working title — revisit once the specific target channel (Sec.~III) is chosen.

\author{Author list TBD}
\affiliation{Affiliation TBD}

\date{\today}

\begin{abstract}
Reference~\cite{perruzza2026} proposed a pair of Coulomb-entangled electrons on superfluid helium,
prepared in the motional/spin singlet, as a differential sensor for particle-physics signals,
demonstrating a factor-of-two quantum-Fisher-information gain over an independent-spin Ramsey
benchmark for a \emph{static} magnetic-field gradient. That work explicitly deferred time-dependent
fields to future study. This project takes up that extension: we generalize the singlet--triplet
sensor Hamiltonian to an oscillating or transient coupling $\lambda(t)$ and ask which particle-physics
targets are best matched to it. A literature survey of running and proposed experiments —
oscillating axion/axion-like-particle (ALP) dark-matter backgrounds (CASPEr, ABRACADABRA, DM~Radio),
dark-photon dark matter, network-based domain-wall/transient searches (GNOME), and precision
electric-dipole-moment (EDM) measurements — is used to identify which regime (resonant narrowband,
broadband stochastic, or impulsive) the eHe platform's characteristic frequency scale and
common-mode-rejection architecture could plausibly compete in. This is a proposal: we outline the
extended formalism, the candidate target channels with their existing experimental sensitivity
reach, and the open questions that must be resolved before a concrete protocol can be simulated.
\end{abstract}

\maketitle

\section{Introduction}
\label{sec:intro}

Two Coulomb-entangled electrons on superfluid helium, prepared in the motional/spin singlet, were
shown in Ref.~\cite{perruzza2026} to act as a differential quantum sensor: blind to common-mode
fields and maximally responsive to a \emph{difference} between the two trap sites, quantified through
the quantum and classical Fisher information and benchmarked against an independent-spin Ramsey
scheme. The entangled states and the numerical machinery used to obtain them (Hartree/CI on a
sine-DVR basis, configurations I/II/III) were established in Ref.~\cite{beysengulov2024}, with
time-dependent \emph{control} of the same states realized as high-fidelity two-qubit gates in
Ref.~\cite{leinonen2025}; Ref.~\cite{perruzza2026} also grounds the sensing readout in the real
resonator/charge-readout hardware demonstrated in Refs.~\cite{koolstra2025impedance,castoria2025}.

Reference~\cite{perruzza2026} is explicit about its scope: ``we limit ourselves in this work to the
detection of a static magnetic-field gradient across the two wells,'' and states that ``we will
address time-dependent fields in future studies.'' The same paper's discussion section already
gestures at what those time-dependent targets might be — ``the extremely weak, slowly oscillating
fields associated with ultralight dark-matter candidates such as axions or axion-like
particles''~\cite{graham2013new,abel2017search,budker2014proposal}, entangled sensor networks for
dark-matter and axion signals~\cite{brady2022entangled,derevianko2018detecting}, ultrafast transients
from milli-charged particles or hidden-sector bosons, and rare-event searches including electric
dipole moments — without developing any of them quantitatively. The paper's own sensor-network
outlook (Sec.~2.9) further distinguishes two coupling channels that a single eHe platform could in
principle host: the singlet, optimal for \emph{gradient} signals (generator $G_-$), and Bell/GHZ-type
states, optimal for \emph{common-mode} oscillating signals (generator $G_+$) such as ``a weak
oscillating field from an axion-like background.''

This project asks two questions left open by Ref.~\cite{perruzza2026}: (i) how does the static-field
QFI/CRB formalism of Sec.~IV of that paper generalize to a periodic or transient $\lambda(t)$, and
(ii) among the particle-physics targets it names only qualitatively, which one is quantitatively the
best match for the eHe platform's trap geometry, coherence budget, and $G_\pm$ readout modes, given
where existing dedicated experiments already sit in sensitivity and frequency/mass coverage. We do
not yet have a settled answer to either question; the goal of this manuscript, at this stage, is to
lay out the extended model and a literature-grounded comparison rather than to report a result.

\section{Model}
\label{sec:model}

\subsection{From a static to a time-dependent coupling}

The starting point is the effective two-level (singlet-like/triplet-like) model of
Ref.~\cite{perruzza2026},
\begin{equation}
H_{\rm eff}(t) = E_g|g\rangle\langle g| + E_e|e\rangle\langle e| +
\lambda(t)\left(|g\rangle\langle e| + |e\rangle\langle g|\right),
\end{equation}
where the original proposal fixes $\lambda = (g\mu_B \delta B/2) M_{ge}$ for a \emph{static} field
gradient $\delta B = B_L - B_R$, with matrix element $M_{ge} = \langle e|S^z_L - S^z_R|g\rangle$. The
natural generalization replaces the static $\delta B$ with a time-dependent source term. Two limiting
forms recur throughout the axion/dark-photon haloscope literature and map directly onto candidate
targets identified below:
\begin{itemize}
  \item \textbf{Narrowband/resonant:} $\delta B(t) = \delta B_0\cos(\omega_a t + \phi)$, appropriate
  for a coherent ultralight bosonic dark-matter field of mass $m_a = \hbar\omega_a$ coherent over the
  local dark-matter coherence time $\tau_c \sim 10^6/\omega_a$ — the regime targeted by
  CASPEr~\cite{budker2014proposal}, ABRACADABRA~\cite{ouellet2019design,salemi2021search}, and
  DM~Radio~\cite{silvafeaver2017design}.
  \item \textbf{Impulsive/transient:} $\delta B(t) = \delta B_0\, f(t/\tau)$ for a short pulse of
  duration $\tau$ crossing the trap, appropriate for a milli-charged particle, hidden-sector boson, or
  domain-wall crossing — the regime targeted by network searches such as
  GNOME~\cite{afach2021search,afach2023gnome}.
\end{itemize}
In both cases the quantum and classical Fisher information of Sec.~IV of Ref.~\cite{perruzza2026}
must be re-derived for a coupling with explicit time dependence, replacing the single scalar
parameter $\delta B$ with either an amplitude-and-phase pair $(\delta B_0,\phi)$ at fixed known
$\omega_a$ (matched-filtering regime, as in resonant haloscopes) or a fully unknown waveform (change-
point/transient-detection regime, as in network searches). Neither extension is a straightforward
relabeling of the existing static-parameter QFI/CRB result; both require revisiting which quantity is
actually being estimated.

\subsection{Common-mode vs. gradient coupling}

Whether a given target couples through the gradient generator $G_-$ (as in the original static
proposal) or the common-mode generator $G_+$ depends on the physical origin of the signal. An
oscillating axion-induced pseudo-magnetic field sourced by the galactic dark-matter background is, to
leading order, spatially uniform on trap length scales ($\sim\mu$m), favoring a $G_+$ (Bell/GHZ)
readout as anticipated in Ref.~\cite{perruzza2026}'s outlook; a signal from a particle traversing (or
passing near) one of the two wells, or a topological-defect/domain-wall crossing with a finite
spatial extent, instead favors the $G_-$ (singlet, gradient) readout already developed for the static
case. Dark-photon dark matter can source either an effectively common-mode oscillating field (through
kinetic mixing with the photon, felt uniformly by both wells) or, if the dark-photon Compton
wavelength becomes comparable to the trap baseline, a genuine gradient signal — a distinction not
present in the original static proposal and requiring its own treatment.

\section{Candidate target channels (to be evaluated)}
\label{sec:channels}

\begin{itemize}
  \item \textbf{Oscillating axion/ALP dark-matter background.} The best-motivated extension of
  Ref.~\cite{perruzza2026}'s own framing: axion/ALP dark matter couples to electron (and nucleon)
  spins as an effective, slowly oscillating Zeeman-like term~\cite{graham2013new,abel2017search}.
  Dedicated experiments (CASPEr-Electric/Gradient~\cite{budker2014proposal}, the inductively coupled
  ABRACADABRA-10\,cm~\cite{ouellet2019design,salemi2021search}, and the lumped-element DM~Radio
  pathfinder~\cite{silvafeaver2017design}) already cover a broad low-mass window
  ($\sim 10^{-14}$--$10^{-6}\,{\rm eV}$); recent work also shows that the same NMR/comagnetometer-
  style apparatus is simultaneously sensitive to the axion-photon coupling and to kinetically mixed
  dark photons~\cite{beadle2025nmr}. The open question is whether the eHe platform's trap-baseline
  and coherence-time scale places it in a genuinely uncovered mass window, or whether it would only
  reproduce existing haloscope sensitivity with a different (entanglement-enhanced) prefactor.
  \item \textbf{Dark-photon dark matter.} A kinetically mixed dark photon sources an oscillating
  effective electric or magnetic field inside a shielded volume, detectable with the same
  spin-precession or charge-coupled apparatus used for axions, with a
  distinguishable frequency/coupling signature~\cite{beadle2025nmr}. Since the eHe platform already
  has both a spin ($G_-$/$G_+$ Zeeman) and a charge ($n_L - n_R$, cf. the
  \texttt{majorana-two-electron-sensor} project's charge channel) coupling route, this may be the
  channel best suited to exploiting the platform's dual readout.
  \item \textbf{Network transient/domain-wall searches.} GNOME-style
  searches~\cite{afach2021search,afach2023gnome} correlate transient signals across a global network
  of comagnetometers to detect topological-defect dark matter crossing the Earth. An array of
  entangled eHe sensors (cf. Ref.~\cite{perruzza2026} Sec.~2.9 and the
  \texttt{rl-agentic-quantum-control} project's optimization machinery) could in principle contribute
  a single, high-coherence node to such a network, or motivate an analogous laboratory-scale
  multi-trap correlation search; the impulsive-signal QFI/CRB treatment above is the theoretical
  prerequisite.
  \item \textbf{Electric dipole moment (EDM) precision channel.} A static electric field applied
  across the trap, combined with a comagnetometer-style oscillating readout, is the standard EDM
  protocol~\cite{aoki2021quantum}; entangled-sensor EDM proposals in trapped-atom systems report
  projected sensitivity below $10^{-30}\,e\cdot{\rm cm}$. This channel is intermediate between the
  original static-gradient proposal (a static source field) and the oscillating targets above (an
  oscillating readout is still required to reject systematics), and would need its own assessment of
  whether the eHe platform's achievable $E$-field and coherence time are competitive.
  \item \textbf{Broader spin-dependent short-range forces.} Reviewed in
  Ref.~\cite{jiang2024searches}, ``fifth-force'' searches for exotic spin-dependent interactions are a
  natural extension of the gradient ($G_-$) channel already built for Ref.~\cite{perruzza2026}, since
  they are inherently short-range, differential signals rather than galactic-background common-mode
  ones — potentially the lowest-effort extension of the existing static formalism, trading a
  particle-physics dark-matter motivation for a table-top BSM-force one.
\end{itemize}

Each channel implies a different generator ($G_-$ vs.\ $G_+$), a different estimation-theoretic
regime (resonant narrowband vs.\ transient), and a different existing-experiment baseline to beat or
complement; narrowing this list — using, at minimum, an order-of-magnitude comparison between the
eHe platform's coherence-time-limited sensitivity and the published exclusion curves of
Refs.~\cite{budker2014proposal,ouellet2019design,salemi2021search,silvafeaver2017design,afach2021search,afach2023gnome,beadle2025nmr}
— is the first concrete task for this project. A broader Snowmass-style survey of quantum sensors for
high-energy physics~\cite{cecil2022snowmass} is used as a cross-check that no obviously more
appropriate existing platform is being overlooked.

\section{Open theoretical questions}
\label{sec:open}

\begin{itemize}
  \item How does the quantum Cram\'er--Rao bound of Ref.~\cite{perruzza2026} generalize when the
  parameter of interest is an oscillation amplitude at a known (or partially known) frequency, versus
  an unknown transient waveform? Is a Bayesian/matched-filter treatment, as used in the haloscope
  literature, more natural here than the single-shot static-parameter QFI used previously?
  \item Does the factor-of-two entanglement advantage found for the static gradient survive, strengthen,
  or degrade once averaged over a stochastic (finite dark-matter-coherence-time) oscillating source,
  as opposed to a fixed classical field?
  \item What eHe trap coherence time and baseline (electron--electron separation) would be needed to
  reach a mass/frequency window not already covered by CASPEr, ABRACADABRA, DM~Radio, or GNOME, and is
  that window physically achievable given the device parameters used in
  Ref.~\cite{beysengulov2024,perruzza2026}?
  \item For the network/transient channel, does the existing single-pair sensor need to be embedded in
  a multi-trap array (per Ref.~\cite{perruzza2026} Sec.~2.9) before a meaningful comparison to GNOME-
  style searches is even possible, or is a single high-coherence node already a useful complementary
  data point?
\end{itemize}

\section{Results}

Not yet available — this section will report the channel-comparison order-of-magnitude estimates and,
for the strongest candidate, a time-dependent QFI/CRB derivation analogous to Sec.~IV of
Ref.~\cite{perruzza2026}.

\section{Discussion and outlook}

If a viable channel is identified, next steps are: (i) extend the CI/Hartree machinery of
Ref.~\cite{beysengulov2024} and the numerical toolkit described in the project's
\texttt{ci-dvr-hartree-numerics} skill to propagate the entangled state under an oscillating or
transient $\lambda(t)$; (ii) derive the corresponding time-dependent quantum and classical Fisher
information, including the effect of finite dark-matter coherence time where relevant; (iii) compare
the achievable sensitivity, as a function of trap baseline and coherence time, against the published
exclusion curves for the channel chosen; (iv) assess whether RL-based trap optimization (see the
companion \texttt{rl-agentic-quantum-control} project) can be re-targeted from the static-gradient
reward used in Ref.~\cite{perruzza2026} to an oscillating-signal reward. This project is
complementary to, not a replacement for, the \texttt{majorana-two-electron-sensor} project, which
explores a condensed-matter (rather than particle-physics) target for the same differential-sensor
machinery.

\begin{acknowledgments}
Placeholder.
\end{acknowledgments}

\bibliography{references}

\end{document}
